\vspace{-0.5cm}
\section{Conclusion}
\label{sec:conclusion}
\vspace{-0.4cm}
% \haodong{rewrite this paragraph according to the abs.}

% By introducing the concept of the effective reuse rate, $E(K, N)$, we show that multi-epoch training leads to a significant reuse of data, with the effective dataset size growing logarithmically with $N$ and saturating to $K$ for large $N$. For power-law Hessian spectrum distributions, this scaling behavior holds, albeit at a different rate before saturation. Our theoretical findings align with recent empirical studies, supporting the idea that multi-epoch training (up to four epochs) leads to negligible changes in test loss compared to training on fresh data at each epoch. These insights highlight the importance of considering both dataset size and distribution when investigating the scaling laws of LLM training and provide a foundation for more accurate modeling of data reuse in future research.


In this paper, we intrdouce the concept of effective reuse rate $E(K,N)$ and show both theoretically and empirically that larger datasets can be repeated more. More specifically, $E(K, N)$ saturates at later epochs as the dataset size increases.
% \section{Conclusion}
%In this paper, we characterize how multi-epoch training reshapes data scaling laws through the notion of effective reuse rate $E(K,N)$, defined as the multiplicative factor by which the dataset must grow under one-pass training to achieve the same test loss as $K$-epoch training on $N$ samples. In linear regression with SGD, we prove that when $K$ is small, $E(K,N)\approx K$; as $K$ grows, $E(K,N)$ plateaus at a value of order $\Theta(\log N)$ under strong convexity and at a power of $N$ under a data distribution with power-law structure. Therefore, repeating data is not equivalent to scaling by a constant independent of $N$; larger datasets can be repeated more before returns diminish. 
Several directions remain open for future study. 
(i)~Our analysis is limited to the setting of optimal constant learning rates, and it remains to understand how learning rate decay, as well as the effort spent on tuning the learning rate, affect the effective reuse rate.
(ii)~Our theoretical analysis currently focuses on linear models, and a natural next step would be to extend the framework to more complex settings, such as neural networks with feature learning.
(iii)~We focus on reusing the entire dataset across multiple epochs, but to fully explore the potential of data reuse, one may consider more refined strategies such as selectively reusing or rephrasing high-quality data.
(iv)~Technically, our main results rely on strong convexity. Although we analyze a solvable case under a Zipf-law data distribution, generalizing these proof ideas to broader non-strongly convex settings remains an important theoretical challenge.


%to fully explore the potential of data reuse, one may consider more refined strategies such as data mixing, curriculum learning, or selectively reusing high-quality data.

%(i) Our analysis is limited to the setting of using optimal constant learning rates, and it would be interesting to study how learning rate decay and how the efforts for tuning the learning rate affect the effective reuse rate. (ii) Our analysis is limited to the linear model, and it would be interesting to extend the framework to more complex and realistic settings, such as neural networks with feature learning. (iii) Our work focuses on reusing the whole dataset with multiple epochs. However, to fully explore the potential of data reuse, one can consider a more efficient and heuristic approach to repeating data, such as data mixing, curriculum learning, or reusing only high-quality data.
%(iv) Technically, our main results rely on strong convexity. In the non-strongly convex regime, we provide a solvable case with a Zipf-law data distribution. It would be interesting to generalize these proof ideas to general non-strongly convex linear regression.

%\kaifeng{we should mention the discussion on learning rate decay!}

\section*{Acknowledgements}
This work is supported by the National Natural Science Foundation of China under Grant Number 62495062. Binghui Li is supported by the Elite Ph.D. Program in Applied Mathematics at Peking University. We would like to thank Licong Lin, Huaqing Zhang, Jingfeng Wu, Kaiyue Wen, Jingzhao Zhang and Jason D. Lee for their insightful comments, and the
anonymous reviewers for their valuable feedback.
